\begin{theorem}[classification in Viaclovsky's coordinates]\hypertarget{tag-032S}{}\label{push:src:E.1.2}
\StatementRecord{push:src:E.1.2}{proved}{032S}{\SeriesRef{push:src:E.1.1}}

Use a local logarithm $c$ of the transported linearization,
$u=e^{-2\pi i c}$, and the same reference section and cusp comparison
on both sides.  Two sphere members with data
$(m,\rho_0,N_0;n,\rho_1,N_1;c)$ and primed data are biholomorphic
if and only if their multiplicities agree and there are
$\alpha\in\{1,-1\}$ and $h=c'-c\in\mathbb Z$ such that
\[
 b'=\alpha b,\qquad
 hN'_0\equiv\rho'_0-\alpha\rho_0\pmod{k},\qquad
 hN'_1\equiv\rho'_1-\alpha\rho_1\pmod{2l}.
\]
Every such biholomorphism covers the identity of $B$.
For fixed discrete data the period identification subgroup is
$2kl\mathbb Z$.  At fixed $u$ all sphere translation data give
exactly $2kl$ classes.  A positive $u$-circuit acts on them by a
single cycle.
\end{theorem}


\begin{proof}
\begin{enumerate}[label=\arabic*.,leftmargin=*]
\item We first give the complete sufficiency argument in the stated problem's
clockwise coordinates, writing $w_j=r_j\varepsilon_j+s_j\widehat\delta$.
Set $S_h=I+hE$ and $d_j=S_hw'_j-\alpha w_j$.
The repeated norm maps have images
\[
 N_m\Lambda=m\mathbb Z\varepsilon+k\mathbb Z\widehat\delta,
 \qquad
 N_n\Lambda=n\mathbb Z\varepsilon'+2l\mathbb Z\widehat\delta.
\]
This follows directly from
\[
 N_3v=3\gamma(v)\varepsilon+(2v_u+v_w+3v_\delta)\widehat\delta,
 \quad
 N_4v=4\gamma(v)\varepsilon'+(2v_u+2v_w+4v_\delta)\widehat\delta.
\]
Suppose $p'=\alpha p$ and
$hr'_1\equiv\alpha s_1-s'_1\pmod k$,
$hr'_2\equiv\alpha s_2-s'_2\pmod{2l}$.
Reduction of $p'=\alpha p$ modulo $m,n$ gives
$r'_j\equiv\alpha r_j\pmod{m_j}$.
We may choose integral solutions of $N_{m_j}\ell_j=-d_j$ explicitly:
\[
 \begin{split}
 \ell_1&=\left(\frac{\alpha r_1-r'_1}{m},0,
 -\frac{s'_1+hr'_1-\alpha s_1}{k},0\right)^t,\\
 \ell_2&=\left(\frac{\alpha r_2-r'_2}{n},
 -\frac{s'_2+hr'_2-\alpha s_2}{2l},0,0\right)^t.
 \end{split}
\]
Put $y_j=d_j/m_j+\ell_j$; these vectors belong to the rational
kernel of $N_{m_j}$ and have first coordinate zero.
The following rational vectors solve $(I-A_j)b_j=y_j$:
\[
 \begin{split}
 b_1&=(0,(2y_{1u}+y_{1w})/3,(y_{1w}-y_{1u})/3,0)^t,\\
 b_2&=(0,(y_{2u}-y_{2w})/2,(y_{2u}+y_{2w})/2,0)^t.
 \end{split}
\]
The last coordinates of the equations follow from the displayed
norm-kernel equations.  The maps
\[
 (s,\zeta)\longmapsto(s,\alpha\zeta+\Pi_c(s)b_j)
\]
are holomorphic maps of the normalized torus families over the
degree-$m_j$ root discs.  Their inverses are holomorphic, and period
covariance and $(I-A_j)b_j=d_j/m_j+\ell_j$ give exact commutation
with the prescribed affine deck generators, modulo the integral
period $\ell_j$.  They therefore descend to biholomorphisms of the
finite fillings, including the full multiple fibres.

\item On the punctured root collar the corresponding middle translation is
\[
 A_j=\Pi_cb_j+\frac{\log s_j}{2\pi i}\Pi_cd_j.
\]
Its integral defect is $\ell_j$.  The middle cocycle relation fixes
$\ell_0=-A_0(\ell_1+A_1\ell_2)$.  In the quotient
$\overline\Lambda=\Lambda/(\mathbb Z\widehat w+\mathbb Z\widehat\delta)$,
the first-coordinate condition is exactly $p'=\alpha p$.
Adding an integral multiple of $(0,1,-2,0)^t\in\ker N_m$ to
$\ell_1$ changes the other quotient coordinate by a generator,
without changing this first coordinate.  We can consequently arrange
\[
 [\ell_0]=-\alpha c(\sigma)+c(\sigma')\quad\hbox{in }\overline\Lambda.
\]
An integral toric vector then supplies $q\in\mathbb Z^2$ such that
the cusp reference $sq+\alpha\sigma-\sigma'$ has precisely defect
$\ell_0$.  Here $s=(2\pi i)^{-1}\log t_c$ is the cusp logarithm;
the maps of cusp toric covers are
\[
 (x,t_c)\longmapsto(t_c^q a(t_c)x^\alpha,t_c),
 \qquad a(0)\in(\mathbb C^*)^2.
\]
Integer translations of the cocharacter lattice, its inversion,
and multiplication by a holomorphic unit preserve the specified
hexagonal fan and all its strata.  These maps and their inverses
extend holomorphically across the three double curves and both
triple points.

\item We justify the global analytic correction.  Differences of allowed
finite extensions are the equivariant holomorphic vertical vectors
on the complete root discs.  The extra subgroup of order $k$ at
the first point acts trivially on the vertical vector coefficients
and rotates $s$; its invariant holomorphic coefficients are exactly
those in $s^k$.  The remaining order-three equivariance is the
original one.  The argument at the second point uses $s^l$ and
order four.  Thus the coarse difference sheaf is literally the same
sheaf $\mathcal V_B$ as for multiplicities $(3,4)$, not an assumed
analogue.  Its first-coordinate quotient is
$\mathcal O_B(-p_0)$: the modular frame has vanishing orders $2k,l$
on the two enlarged roots and a simple cusp pole.  Its fixed
second-coordinate subline is $\mathcal O_B$.  Finite averaging
provides local lifts, hence
\[
 0\longrightarrow\mathcal O_B\longrightarrow\mathcal V_B
 \longrightarrow\mathcal O_B(-p_0)\longrightarrow0.
\]
Both end terms have vanishing $H^1$ on $\mathbb P^1$, so
$H^1(B,\mathcal V_B)=0$.  The permitted local translations form a
torsor under this sheaf: on each overlap their defects agree, so
their difference is a single-valued section of the indicated
difference sheaf.  Vanishing of its \v Cech class gives a global
middle translation.  Its overlap equalities are exactly the two
finite formulas and the cusp formula just displayed.  Consequently
the local biholomorphisms glue, with their inverses, over $B$.

\item For necessity the base is intrinsic even when its anticanonical
pencil changes.  If $F:X\to X'$ is a biholomorphism, the pullback
of the Fubini--Study form by $f'F$ is exact because
$H^2(S^6,\mathbb R)=0$.  Its restriction to a regular complex torus,
wedged with a flat K\"ahler form, has zero integral.  Positivity
forces $f'F$ to be constant on that torus.  The proper analytic
image of $(f,f'F)$ is therefore a curve of degree one over $B$;
normalization and the inverse argument give a base automorphism.
Its three monodromy orders $3,4,\infty$ are distinct, so it fixes
all three marked points and is the identity.  Multiplicities agree.
On ordinary fibres the lift is affine.  The simultaneous rational
centralizer of $A_1,A_2$ is $\mathbb Q I+\mathbb Q E$, as follows
by solving their entrywise commuting equations.  Comparison of the
constant last two columns of $\Pi_c$ and unimodularity then gives
fibre-linear part $\alpha I$, integer difference $h$, and lattice
map $\alpha S_h^{-1}$.  The normalized root base change is
functorial, so the lifted map commutes with the actual deck
generator.  Summing its defect around that generator gives the
two norm conditions above.  A cusp map has an integral toric
winding: at a generic point of a toric component, its ratio of
characters is $t_c^q$ times a holomorphic unit.  Its class in
$\overline\Lambda$ vanishes.  The middle cocycle relation then
gives $p'=\alpha p$.  This proves all the necessary conditions.

\item Substitute \ref{push:src:E.1.1}.  In the second congruence the extra term
$N'_1-\alpha N_1$ is divisible by $4l$, and hence drops out
modulo $2l$.  We obtain exactly the statement's two congruences.
For fixed data $b\ne0$ excludes $\alpha=-1$.
The invertibility of $N_0$ modulo $k$, of $N_1$ modulo $2l$, and
$\gcd(k,2l)=1$ give $h\in2kl\mathbb Z$.
For varying data the Chinese remainder theorem gives exactly one
shift class modulo $2kl$, and all residue classes of
$\rho_0\bmod k$, $\rho_1\bmod2l$ occur without changing $b$ or
freeness.  Finally a positive $u$-circuit sends $c$ to $c-1$.
Its translation on $\mathbb Z/(2kl)$ is a generator.
\end{enumerate}
\end{proof}